% GENERATED FILE. Edit canonical lessons, then run npm run generate:books.
% canonical-source-hash: fnv1a64:35864590418b5c3a
% canonical-lessons: GE-C397-abstuerzen, GE-C397-aktualisieren, GE-C397-fehlermeldung, GE-C397-wiederherstellen, GE-C397-sicherung, GE-R397-review-394-397, GE-R397-review-final-397

\chapter{Der Computer stürzt ab}
\label{ch:ge-der-computer-sturzt-ab}
\hlchaptermodality{397}

\begin{chapteropening}
\textbf{By the end of this chapter:} \emph{I can say that a computer has crashed or needs updating, read out an error message, and say how data was restored from a backup.}

Complete the last lesson of chapter 397: I can say that a computer has crashed or needs updating, read out an error message, and say how data was restored from a backup.
\end{chapteropening}

\section[abstürzen]{abstürzen — to crash; to fall (from a height)}

\label{lesson:GE-C397-abstuerzen}



\begin{quote}
Before the new one: say the German for to install, then the German for the program.
\end{quote}

\subsection*{You'll want to know: abstürzen}
\textbf{abstürzen} — \textquotedblleft{}to crash; to fall (from a height)\textquotedblright{}.

\begin{quote}
Mein Computer ist schon wieder abgestürzt.
\end{quote}

\begin{quote}
\emph{My computer has crashed yet again.}
\end{quote}

Separable, Perfekt with \emph{sein}: \emph{er stürzt ab}, \emph{er ist abgestürzt}. Also of planes and climbers.

\begin{cousinweb}
\emph{abstürzen} is built from \emph{ab} \textquotedblleft{}down, off\textquotedblright{} and \emph{stürzen} \textquotedblleft{}to plunge\textquotedblright{}, inherited from Middle High German \emph{stürzen}, Old High German \emph{sturzen} and Proto-Germanic \emph{sturtijaną} (EtymDB).
\end{cousinweb}

\begin{grammarlens}[title={What you've built}]
One more word for telling, explaining and describing at B1.
\end{grammarlens}

\subsection*{Your turn}
\begin{itemize}
\raggedright
  \item \emph{Say it:} \emph{abstürzen}
  \item \emph{Say it:} \emph{abstürzen}, once more
  \item \emph{Say it:} a sentence of your own with \emph{abstürzen}
  \item \emph{Recall:} say the German for to install, then the German for the program, then say \emph{abstürzen} again
\end{itemize}

\begin{tcolorbox}[breakable,colback=teal!4,colframe=teal!35!black,title={Before you move on}]
What does abstürzen mean? (\textquotedblleft{}To crash; to fall (from a height)\textquotedblright{}.) Say it once more.
\end{tcolorbox}
\section[aktualisieren]{aktualisieren — to update}

\label{lesson:GE-C397-aktualisieren}



\begin{quote}
Before the new one: say the German for the program, then the German for to crash.
\end{quote}

\subsection*{You'll want to know: aktualisieren}
\textbf{aktualisieren} — \textquotedblleft{}to update\textquotedblright{}.

\begin{quote}
Das Handy aktualisiert die Apps in der Nacht.
\end{quote}

\begin{quote}
\emph{The phone updates the apps during the night.}
\end{quote}

Perfekt: \emph{hat aktualisiert}. From \emph{aktuell} — current, up to date; a false friend, not actual.

\begin{cousinweb}
\emph{aktualisieren} is built on \emph{aktuell} \textquotedblleft{}current\textquotedblright{}, from Late Latin \emph{āctuālis} \textquotedblleft{}active, practical\textquotedblright{} (EtymDB; Lewis \& Short), from \emph{āctus} \textquotedblleft{}a doing\textquotedblright{}, of the family of \emph{agere} \textquotedblleft{}to do, to drive\textquotedblright{}.
\end{cousinweb}

\begin{grammarlens}[title={What you've built}]
One more word for telling, explaining and describing at B1.
\end{grammarlens}

\subsection*{Your turn}
\begin{itemize}
\raggedright
  \item \emph{Say it:} \emph{aktualisieren}
  \item \emph{Say it:} \emph{aktualisieren}, once more
  \item \emph{Say it:} a sentence of your own with \emph{aktualisieren}
  \item \emph{Recall:} say the German for the program, then the German for to crash, then say \emph{aktualisieren} again
\end{itemize}

\begin{tcolorbox}[breakable,colback=teal!4,colframe=teal!35!black,title={Before you move on}]
What does aktualisieren mean? (\textquotedblleft{}To update\textquotedblright{}.) Say it once more.
\end{tcolorbox}
\section[die Fehlermeldung]{die Fehlermeldung — the error message}

\label{lesson:GE-C397-fehlermeldung}



\begin{quote}
Before the new one: say the German for to crash, then the German for to update.
\end{quote}

\subsection*{You'll want to know: die Fehlermeldung}
\textbf{die Fehlermeldung} — \textquotedblleft{}the error message\textquotedblright{}.

\begin{quote}
Beim Start kommt jedes Mal eine Fehlermeldung.
\end{quote}

\begin{quote}
\emph{Every time it starts, an error message comes up.}
\end{quote}

Plural: \emph{die Fehlermeldungen}.

\begin{cousinweb}
\emph{Fehlermeldung} joins \emph{Fehler} \textquotedblleft{}mistake\textquotedblright{}, from \emph{fehlen} \textquotedblleft{}to be missing\textquotedblright{}, to \emph{Meldung} \textquotedblleft{}report\textquotedblright{}, from \emph{melden} \textquotedblleft{}to report\textquotedblright{}, Old High German \emph{meldōn} (EtymDB).
\end{cousinweb}

\begin{grammarlens}[title={What you've built}]
One more word for telling, explaining and describing at B1.
\end{grammarlens}

\subsection*{Your turn}
\begin{itemize}
\raggedright
  \item \emph{Say it:} \emph{die Fehlermeldung}
  \item \emph{Say it:} \emph{die Fehlermeldung}, once more
  \item \emph{Say it:} a sentence of your own with \emph{die Fehlermeldung}
  \item \emph{Recall:} say the German for to crash, then the German for to update, then say \emph{die Fehlermeldung} again
\end{itemize}

\begin{tcolorbox}[breakable,colback=teal!4,colframe=teal!35!black,title={Before you move on}]
What does die Fehlermeldung mean? (\textquotedblleft{}The error message\textquotedblright{}.) Say it once more.
\end{tcolorbox}
\section[wiederherstellen]{wiederherstellen — to restore}

\label{lesson:GE-C397-wiederherstellen}



\begin{quote}
Before the new one: say the German for to update, then the German for the error message.
\end{quote}

\subsection*{You'll want to know: wiederherstellen}
\textbf{wiederherstellen} — \textquotedblleft{}to restore\textquotedblright{}.

\begin{quote}
Der Techniker konnte die alten Fotos wiederherstellen.
\end{quote}

\begin{quote}
\emph{The technician was able to restore the old photos.}
\end{quote}

Separable: \emph{er stellt wieder her}, \emph{er hat wiederhergestellt}.

\begin{cousinweb}
\emph{wiederherstellen} is built from \emph{wieder} \textquotedblleft{}again\textquotedblright{} and \emph{herstellen} \textquotedblleft{}to make, to produce\textquotedblright{}, from \emph{stellen} \textquotedblleft{}to put\textquotedblright{}, Proto-Germanic \emph{stallijaną} (EtymDB): you put a thing back as it was.
\end{cousinweb}

\begin{grammarlens}[title={What you've built}]
One more word for telling, explaining and describing at B1.
\end{grammarlens}

\subsection*{Your turn}
\begin{itemize}
\raggedright
  \item \emph{Say it:} \emph{wiederherstellen}
  \item \emph{Say it:} \emph{wiederherstellen}, once more
  \item \emph{Say it:} a sentence of your own with \emph{wiederherstellen}
  \item \emph{Recall:} say the German for to update, then the German for the error message, then say \emph{wiederherstellen} again
\end{itemize}

\begin{tcolorbox}[breakable,colback=teal!4,colframe=teal!35!black,title={Before you move on}]
What does wiederherstellen mean? (\textquotedblleft{}To restore\textquotedblright{}.) Say it once more.
\end{tcolorbox}
\section[die Sicherung]{die Sicherung — the backup; the fuse}

\label{lesson:GE-C397-sicherung}



\begin{quote}
Before the new one: say the German for the error message, then the German for to restore.
\end{quote}

\subsection*{You'll want to know: die Sicherung}
\textbf{die Sicherung} — \textquotedblleft{}the backup; the fuse\textquotedblright{}.

\begin{quote}
Zum Glück hatte ich eine Sicherung von allen Dateien.
\end{quote}

\begin{quote}
\emph{Luckily I had a backup of all the files.}
\end{quote}

Plural: \emph{die Sicherungen}. Also a fuse: \emph{Die Sicherung ist durchgebrannt.} Many people also say \emph{das Backup}.

\begin{cousinweb}
\emph{Sicherung} is the noun of \emph{sichern} \textquotedblleft{}to make safe\textquotedblright{} (EtymDB), from \emph{sicher} \textquotedblleft{}safe\textquotedblright{}, which goes back through Old High German \emph{sihhur} to Latin \emph{sēcūrus} \textquotedblleft{}free from care\textquotedblright{}.
\end{cousinweb}

\begin{grammarlens}[title={What you've built}]
One more word for telling, explaining and describing at B1.
\end{grammarlens}

\subsection*{Your turn}
\begin{itemize}
\raggedright
  \item \emph{Say it:} \emph{die Sicherung}
  \item \emph{Say it:} \emph{die Sicherung}, once more
  \item \emph{Say it:} a sentence of your own with \emph{die Sicherung}
  \item \emph{Recall:} say the German for the error message, then the German for to restore, then say \emph{die Sicherung} again
\end{itemize}

\begin{tcolorbox}[breakable,colback=teal!4,colframe=teal!35!black,title={Before you move on}]
What does die Sicherung mean? (\textquotedblleft{}The backup; the fuse\textquotedblright{}.) Say it once more.
\end{tcolorbox}
\section[(dialogue)]{Bedingungen und Regeln}

\label{lesson:GE-R397-review-394-397}



\begin{quote}
Say the German for to upload and the web page.
\end{quote}

\subsection*{Your turn}
Say these aloud:

\begin{itemize}
\raggedright
  \item to upload, the web page, the search engine, to post, to comment on
  \item the password, to log in, the access, data protection, to encrypt
  \item the keyboard, the printer, to print out, to install, the program
  \item to crash, to update, the error message, to restore, the backup
\end{itemize}

\begin{tcolorbox}[breakable,colback=teal!4,colframe=teal!35!black,title={Before you move on}]
To upload? (\textbf{hochladen}.) The web page? (\textbf{die Webseite}.) The search engine? (\textbf{die Suchmaschine}.) To post? (\textbf{posten}.) To comment on? (\textbf{kommentieren}.) The password? (\textbf{das Passwort}.) To log in? (\textbf{einloggen}.) The access? (\textbf{der Zugang}.) Data protection? (\textbf{der Datenschutz}.) To encrypt? (\textbf{verschlüsseln}.) The keyboard? (\textbf{die Tastatur}.) The printer? (\textbf{der Drucker}.) To print out? (\textbf{ausdrucken}.) To install? (\textbf{installieren}.) The program? (\textbf{das Programm}.) To crash? (\textbf{abstürzen}.) To update? (\textbf{aktualisieren}.) The error message? (\textbf{die Fehlermeldung}.) To restore? (\textbf{wiederherstellen}.) The backup? (\textbf{die Sicherung}.)
\end{tcolorbox}
\section[(dialogue)]{Der Computer stürzt ab — once more}

\label{lesson:GE-R397-review-final-397}



\begin{quote}
Say the German for to crash and to update.
\end{quote}

\subsection*{Your turn}
\begin{itemize}
\raggedright
  \item \emph{Say it:} to crash, to update, the error message, to restore, the backup
\end{itemize}

\begin{tcolorbox}[breakable,colback=teal!4,colframe=teal!35!black,title={Before you move on}]
To crash? (\textbf{abstürzen}.) To update? (\textbf{aktualisieren}.) The error message? (\textbf{die Fehlermeldung}.) To restore? (\textbf{wiederherstellen}.) The backup? (\textbf{die Sicherung}.)
\end{tcolorbox}
